% Part: first-order-logic
% Chapter: axiomatic-proofs
% Section: identity

\documentclass[../../../include/open-logic-section]{subfiles}

\begin{document}

\olfileid{fol}{axd}{ide}

\olsection{\usetoken{P}{derivation} met \usetoken{S}{identity}: gelijkheid gebruiken}

Om $\eq$ in !!{derivation}s te kunnen gebruiken, voegen we alleen
nieuwe axiomaschema's toe. Verder verandert er niets: de definitie van
!!{derivation} en $\Proves$ blijft hetzelfde. De nieuwe axioma's zijn
voortaan ook toegestaan.

\begin{defn}[Axioma's voor !!{identity}, dus gelijkheid]
\begin{align}
\ollabel{ax:id1} & \eq[t][t], \\
\ollabel{ax:id2} & \eq[t_1][t_2] \lif (!B(t_1) \lif
  !B(t_2)),
\end{align}
Hierbij zijn $t$, $t_1$, $t_2$ willekeurige gesloten termen: termen
zonder vrije variabelen.
\end{defn}

\begin{prop}\ollabel{prop:sound}
  De axioma's \olref{ax:id1} en \olref{ax:id2} zijn geldig.
\end{prop}

\begin{proof}
  Dit bewijs is een oefening.
\end{proof}

\begin{prob}
Bewijs \olref[fol][axd][ide]{prop:sound}.
\end{prob}

\begin{prop}\ollabel{prop:iden1}
 $\Gamma \Proves \eq[t][t]$, voor elke term $t$ en elke
 verzameling~$\Gamma$.
\end{prop}

\begin{prop}\ollabel{prop:iden2}
  Stel dat $\Gamma \Proves !A(t_1)$ en $\Gamma \Proves
  \eq[t_1][t_2]$. Dan geldt $\Gamma \Proves !A(t_2)$.
\end{prop}

\begin{proof}
De !!{formula}
\[
(\eq[t_1][t_2] \lif (!A(t_1) \lif !A(t_2)))
\]
is een instantie van~\olref{ax:id2}. Pas nu \MP{} tweemaal toe. Dan
volgt de conclusie.
\end{proof}

\end{document}
